\cotefile{82}
\scanpage{50}
% @ lines 1–6 / text
Mais il y a dans $Q$ un \unsure{autom.} involutif canonique, qui correspond à la symétrie des deux facteurs $X\times X$ dans (6). Il est induit par un antiautomorphisme \inserted{can.} de l'Algèbre $M$ obtenu \inserted{d'abord $\mathrm{loc}^{\mathrm t}$} en choisissant \deleted{\unsure{loc$^{\mathrm t}$ une base}} une \deleted{\unsure{$\omega$}} base $\omega$ de $\det\check V$\inserted{$=\bigwedge^2\check V$}, d'où un isom \deleted{\illegible{1}}
% @ line 7 / display
\[
 V\xrightarrow{\deleted{\sim}\,\alpha}\check V
 \qquad\text{composé de}\quad V\longrightarrow\check V\otimes\det V\xrightarrow{\omega}\check V
 \tag*{(13)}
\]
% @ line 8 / text
ce qui permet de définir
% @ line 9 / display
\[
 u\longmapsto u'=\alpha^{-1}({}^{t}u)\alpha\tag*{(14)}
\]

% @ lines 10–12 / text
Mais changer de $\omega$ revient à multiplier $\alpha$ par une \deleted{section} unité $\lambda\inG_m(S)$, donc \deleted{\unsure{$u$ ch.}}\inserted{$\alpha^{-1}$ devient $\lambda^{-1}\alpha$} et $u'$ ne change pas. Par construction
% @ line 13 / display
\[
 (uv)'=v'u' \tag*{(15)}
\]

% @ line 14 / text
En fait, le calcul donne
% @ line 15 / display
\[
 u'=\operatorname{Tr}u.1-u\tag*{(16)}
\]
% @ line 16 / text
plus richement on a
% @ lines 17–18 / display
\[
 \left\{\begin{aligned}u+u'&=(\operatorname{Tr}u)\,1\\ uu'&=(\det u)\,1\end{aligned}\right. \quad,\tag*{(17)}
\]
% @ lines 19–21 / text
et il est immédiat que $u\longmapsto u'$ échange $\operatorname{Ker}$ et $\operatorname{Im}$ pour un $u$ tel que $P(u)=0$, \unsure{ce|$u$} non nul sur les fibres....

% @ lines 22–25 / text
Quand on voit ce que devient cette involution de $M$ et $Q$ en tenant compte de (12), on voit qu'à tt couple $(A,\rho)$ elle associe $A$ ($=A'$ !) et $\rho'$ défini par
% @ line 26 / display
\[
 \rho'(u)=\rho(u')\tag*{(18)}
\]
% @ line 26, lower right / margin note
\margin{(\uline{NB} $u'|A$ est un \unsure{autom.|automorphisme}, car $A$ comm.)}
% archive sheet number: 49
